\lecture{07}{Следствия сильной выпуклости}

\sourcevideo
    {https://www.youtube.com/playlist?list=PLOzRYVm0a65fhKDS8cYIC7YCuVGJ4FOJx}
    {Week 2: Lecture 7: Implications of Strong Convexity}

Сильная выпуклость допускает несколько эквивалентных
характеризаций. Одни из них удобны для проверки свойства функции,
другие используются при анализе сходимости оптимизационных
алгоритмов.

\section{Эквивалентные характеризации}

Пусть \(C\subseteq\RR^n\) открыто и выпукло, \(f\in C^1(C)\), а
\(\mu>0\).

\begin{theorem}[Характеризации сильной выпуклости]
Следующие условия эквивалентны.

Функция \(f\) является \(\mu\)-сильно выпуклой, то есть
\[
    f(y)
    \ge
    f(x)
    +
    \inner{\nabla f(x)}{y-x}
    +
    \frac{\mu}{2}\norm{y-x}^2
\]
для всех \(x,y\in C\).

Функция
\[
    g(x)
    =
    f(x)-\frac{\mu}{2}\norm{x}^2
\]
выпукла.

Градиент \(f\) является \(\mu\)-сильно монотонным:
\[
    \inner{
        \nabla f(x)-\nabla f(y)
    }{
        x-y
    }
    \ge
    \mu\norm{x-y}^2.
\]

Для любых \(x,y\in C\) и \(\lambda\in[0,1]\)
\[
\begin{aligned}
    f\bigl(\lambda x+(1-\lambda)y\bigr)
    \le{}&
    \lambda f(x)+(1-\lambda)f(y)
    \\
    &-
    \frac{\mu}{2}
    \lambda(1-\lambda)\norm{x-y}^2.
\end{aligned}
\]
\end{theorem}

\begin{proof}
Поскольку
\[
    \nabla g(x)=\nabla f(x)-\mu x,
\]
критерий выпуклости первого порядка для \(g\) имеет вид
\[
    g(y)
    \ge
    g(x)
    +
    \inner{\nabla f(x)-\mu x}{y-x}.
\]
Подставляя определение \(g\) и собирая квадратичные слагаемые,
получаем
\[
    f(y)
    \ge
    f(x)
    +
    \inner{\nabla f(x)}{y-x}
    +
    \frac{\mu}{2}\norm{y-x}^2.
\]
Следовательно, первые два условия эквивалентны.

Если \(g\) выпукла, то её градиент монотонен:
\[
    \inner{\nabla g(x)-\nabla g(y)}{x-y}\ge0.
\]
После подстановки
\[
    \nabla g(x)-\nabla g(y)
    =
    \nabla f(x)-\nabla f(y)-\mu(x-y)
\]
получаем третье условие.

Обратно, третье условие означает монотонность \(\nabla g\). Для
\(d=y-x\) имеем
\[
\begin{aligned}
    g(y)-g(x)-\inner{\nabla g(x)}{d}
    &=
    \int_0^1
    \inner{
        \nabla g(x+td)-\nabla g(x)
    }{d}
    \,\mathrm dt
    \\
    &\ge0,
\end{aligned}
\]
поскольку при \(t>0\) монотонность градиента даёт
\[
    \inner{
        \nabla g(x+td)-\nabla g(x)
    }{td}
    \ge0.
\]
Следовательно, \(g\) выпукла.

Наконец, выпуклость \(g\) равносильна неравенству
\[
    g\bigl(\lambda x+(1-\lambda)y\bigr)
    \le
    \lambda g(x)+(1-\lambda)g(y).
\]
Используя тождество
\[
\begin{aligned}
    \norm{\lambda x+(1-\lambda)y}^2
    ={}&
    \lambda\norm{x}^2+(1-\lambda)\norm{y}^2
    \\
    &-
    \lambda(1-\lambda)\norm{x-y}^2,
\end{aligned}
\]
получаем четвёртую характеризацию.
\end{proof}

Таким образом, сильную выпуклость можно понимать как обычную
выпуклость после вычитания квадратичной функции
\[
    \frac{\mu}{2}\norm{x}^2.
\]

\section{Нижняя оценка изменения градиента}

Из сильной монотонности и неравенства Коши--Буняковского следует
\[
\begin{aligned}
    \mu\norm{x-y}^2
    &\le
    \inner{
        \nabla f(x)-\nabla f(y)
    }{
        x-y
    }
    \\
    &\le
    \norm{\nabla f(x)-\nabla f(y)}
    \norm{x-y}.
\end{aligned}
\]
Поэтому при \(x\ne y\)
\[
    \norm{\nabla f(x)-\nabla f(y)}
    \ge
    \mu\norm{x-y}.
\]

В частности, если \(x^\star\) является минимумом и
\(\nabla f(x^\star)=0\), то
\[
    \norm{\nabla f(x)}
    \ge
    \mu\norm{x-x^\star}.
\]
Следовательно, градиент сильно выпуклой функции количественно
контролирует расстояние до оптимального решения.

\section{Липшицевость и гладкость}

Необходимо различать липшицевость самой функции и липшицевость её
градиента.

\begin{definition}[Липшицева функция]
Функция \(f\colon C\to\RR\) называется \(L\)-липшицевой, если
\[
    \abs{f(x)-f(y)}
    \le
    L\norm{x-y}
\]
для всех \(x,y\in C\).
\end{definition}

\begin{definition}[Гладкая функция]
Дифференцируемая функция \(f\) называется \(L\)-гладкой, если её
градиент \(L\)-липшицев:
\[
    \norm{\nabla f(x)-\nabla f(y)}
    \le
    L\norm{x-y}.
\]
\end{definition}

Если функция одновременно \(\mu\)-сильно выпукла и \(L\)-гладка, то
\[
    \mu\norm{x-y}
    \le
    \norm{\nabla f(x)-\nabla f(y)}
    \le
    L\norm{x-y}.
\]
Следовательно,
\[
    0<\mu\le L.
\]

Для функции
\[
    f(x)=x^2
\]
имеем
\[
    \nabla f(x)=2x,
\]
поэтому
\[
    \abs{\nabla f(x)-\nabla f(y)}
    =
    2\abs{x-y}.
\]
Она \(2\)-сильно выпукла и \(2\)-гладка.

Для дважды дифференцируемой функции сочетание сильной выпуклости и
гладкости выражается через Гессиан:
\[
    \mu I
    \preceq
    \nabla^2f(x)
    \preceq
    LI.
\]

\section{Сумма выпуклых функций}

\begin{theorem}
Пусть \(f\) является \(\mu\)-сильно выпуклой, а \(g\) выпукла.
Тогда функция
\[
    h=f+g
\]
также является \(\mu\)-сильно выпуклой.
\end{theorem}

\begin{proof}
Для \(\lambda\in[0,1]\) сильная выпуклость \(f\) даёт
\[
\begin{aligned}
    f\bigl(\lambda x+(1-\lambda)y\bigr)
    \le{}&
    \lambda f(x)+(1-\lambda)f(y)
    \\
    &-
    \frac{\mu}{2}
    \lambda(1-\lambda)\norm{x-y}^2.
\end{aligned}
\]
По выпуклости \(g\)
\[
    g\bigl(\lambda x+(1-\lambda)y\bigr)
    \le
    \lambda g(x)+(1-\lambda)g(y).
\]
Складывая эти неравенства, получаем
\[
\begin{aligned}
    h\bigl(\lambda x+(1-\lambda)y\bigr)
    \le{}&
    \lambda h(x)+(1-\lambda)h(y)
    \\
    &-
    \frac{\mu}{2}
    \lambda(1-\lambda)\norm{x-y}^2.
\end{aligned}
\]
\end{proof}

Следовательно, к выпуклой целевой функции можно добавить сильно
выпуклый регуляризатор и получить сильно выпуклую задачу. Например,
если \(g\) выпукла, то
\[
    h(x)=g(x)+\frac{\mu}{2}\norm{x}^2
\]
является \(\mu\)-сильно выпуклой.

При этом добавление регуляризатора в общем случае изменяет точку
минимума. Сохранение исходного решения требует отдельного
обоснования.

\section{Сильная выпуклость и PL-неравенство}

\begin{theorem}
Пусть \(f\in C^1(\RR^n)\) является \(\mu\)-сильно выпуклой и
\[
    f^\star=\min_x f(x).
\]
Тогда
\[
    \frac12\norm{\nabla f(x)}^2
    \ge
    \mu\bigl(f(x)-f^\star\bigr)
\]
для всех \(x\in\RR^n\).
\end{theorem}

\begin{proof}
Из сильной выпуклости для любого \(y\) следует
\[
    f(y)
    \ge
    f(x)
    +
    \inner{\nabla f(x)}{y-x}
    +
    \frac{\mu}{2}\norm{y-x}^2.
\]
Правая часть как функция \(y\) минимальна при
\[
    y
    =
    x-\frac{1}{\mu}\nabla f(x).
\]
Её минимальное значение равно
\[
    f(x)
    -
    \frac{1}{2\mu}\norm{\nabla f(x)}^2.
\]
Поскольку неравенство справедливо для всех \(y\),
\[
    f^\star
    \ge
    f(x)
    -
    \frac{1}{2\mu}\norm{\nabla f(x)}^2.
\]
После перестановки слагаемых получаем
\[
    \frac12\norm{\nabla f(x)}^2
    \ge
    \mu\bigl(f(x)-f^\star\bigr).
\]
\end{proof}

Таким образом,
\[
    \text{сильная выпуклость}
    \quad\Longrightarrow\quad
    \text{PL-неравенство}.
\]
Обратное утверждение неверно. PL-неравенство может выполняться для
функций, не являющихся сильно выпуклыми и даже выпуклыми.

В лекции вновь упоминается пример
\[
    f(x)=x^2+3\sin^2x.
\]
Он используется для иллюстрации невыпуклой функции, допускающей
PL-анализ. Подробная проверка соответствующего неравенства будет
рассматриваться позднее.

\section{Условие оптимальности на выпуклом множестве}

Рассмотрим задачу
\[
    \minimize_{x\in\mathcal X}f(x),
\]
где \(\mathcal X\subseteq\RR^n\) выпукло, а \(f\) выпукла и
дифференцируема.

\begin{theorem}[Условие первого порядка]
Точка \(x^\star\in\mathcal X\) является глобальным минимумом тогда и
только тогда, когда
\[
    \inner{\nabla f(x^\star)}{y-x^\star}
    \ge0
\]
для всех \(y\in\mathcal X\).
\end{theorem}

\begin{proof}
Пусть \(x^\star\) является минимумом. Для любого \(y\in\mathcal X\)
и \(t\in[0,1]\) точка
\[
    x^\star+t(y-x^\star)
\]
принадлежит \(\mathcal X\). Функция
\[
    \varphi(t)
    =
    f\bigl(x^\star+t(y-x^\star)\bigr)
\]
имеет минимум при \(t=0\), поэтому
\[
    \varphi'(0)
    =
    \inner{\nabla f(x^\star)}{y-x^\star}
    \ge0.
\]

Обратно, по выпуклости
\[
    f(y)
    \ge
    f(x^\star)
    +
    \inner{\nabla f(x^\star)}{y-x^\star}.
\]
Если скалярное произведение неотрицательно, то
\[
    f(y)\ge f(x^\star)
\]
для всех \(y\in\mathcal X\).
\end{proof}

Градиент \(\nabla f(x^\star)\) указывает направление возрастания
функции. Условие оптимальности требует, чтобы угол между градиентом и каждым
допустимым направлением \(y-x^\star\) не превышал \(\pi/2\).

Если минимум лежит на границе допустимого множества, градиент может
быть ненулевым. Он должен быть направлен так, чтобы движение внутрь
множества не уменьшало функцию.

\section{Задача без ограничений}

Если
\[
    \mathcal X=\RR^n,
\]
то условие оптимальности принимает вид
\[
    \inner{\nabla f(x^\star)}{y-x^\star}\ge0
\]
для всех \(y\in\RR^n\).

Можно выбрать
\[
    y=x^\star-\nabla f(x^\star).
\]
Тогда
\[
    -\norm{\nabla f(x^\star)}^2\ge0,
\]
откуда
\[
    \nabla f(x^\star)=0.
\]

Следовательно, для дифференцируемой выпуклой функции
\[
    x^\star\in\argmin_x f(x)
    \quad\Longleftrightarrow\quad
    \nabla f(x^\star)=0.
\]

\section{Квадратичная задача с равенствами}

Рассмотрим задачу
\[
    \begin{aligned}
        \minimize_{x\in\RR^n}
        \quad&
        \frac12x^{\mathsf T}Qx+c^{\mathsf T}x,
        \\
        \subjectto
        \quad&
        Ax=b,
    \end{aligned}
\]
где
\[
    Q=Q^{\mathsf T}\succeq0.
\]
Целевая функция выпукла, а допустимое множество
\[
    \mathcal X=\setgiven{x\in\RR^n}{Ax=b}
\]
аффинно и потому выпукло.

Градиент целевой функции равен
\[
    \nabla f(x)=Qx+c.
\]

По условию первого порядка допустимая точка \(x^\star\) оптимальна
тогда и только тогда, когда
\[
    \inner{Qx^\star+c}{y-x^\star}\ge0
\]
для всех \(y\), удовлетворяющих \(Ay=b\).

Поскольку
\[
    Ay=Ax^\star=b,
\]
имеем
\[
    A(y-x^\star)=0.
\]
Следовательно,
\[
    y-x^\star\in\ker A.
\]

Если \(z\in\ker A\), то допустимы оба направления \(z\) и \(-z\).
Поэтому условие оптимальности эквивалентно равенству
\[
    \inner{Qx^\star+c}{z}=0
    \qquad
    \text{для всех }z\in\ker A.
\]
Иными словами,
\[
    Qx^\star+c
    \in
    (\ker A)^\perp
    =
    \operatorname{range}(A^{\mathsf T}).
\]

В задаче без ограничений это условие сводится к
\[
    Qx^\star+c=0.
\]
Если \(Q\succ0\), то решение единственно и равно
\[
    x^\star=-Q^{-1}c.
\]

При наличии равенства \(Ax=b\) необходимо одновременно выполнить
условие допустимости и условие ортогональности:
\[
    Ax^\star=b,
    \qquad
    Qx^\star+c\perp\ker A.
\]
В следующей лекции эта система будет получена посредством функции
Лагранжа и двойственной задачи.

\begin{lecturesummary}
Сильная выпуклость эквивалентна выпуклости функции
\[
    f(x)-\frac{\mu}{2}\norm{x}^2,
\]
сильной монотонности градиента и усиленному неравенству для выпуклых
комбинаций. Если функция одновременно \(\mu\)-сильно выпукла и
\(L\)-гладка, то
\[
    \mu
    \le
    \frac{
        \norm{\nabla f(x)-\nabla f(y)}
    }{
        \norm{x-y}
    }
    \le
    L.
\]
Сумма \(\mu\)-сильно выпуклой и выпуклой функций остаётся
\(\mu\)-сильно выпуклой. Сильная выпуклость также влечёт
PL-неравенство. Для минимизации выпуклой функции на выпуклом
множестве точка \(x^\star\) оптимальна тогда и только тогда, когда
\[
    \inner{\nabla f(x^\star)}{y-x^\star}\ge0
\]
для всех допустимых \(y\). В задаче без ограничений это условие
переходит в равенство \(\nabla f(x^\star)=0\).
\end{lecturesummary}